% Generated draft: requires mathematical and source review.
\documentclass[11pt]{article}
\usepackage[margin=25mm]{geometry}
\usepackage{fontspec}
\setmainfont{Latin Modern Roman}
\usepackage{amsmath,amssymb,mathtools,unicode-math}
\setmathfont{Latin Modern Math}
\usepackage{enumitem,needspace,xurl,hyperref}
\setlength{\parindent}{0pt}
\setlength{\parskip}{5pt}
\setlength{\emergencystretch}{3em}
\newcommand{\N}{\mathbb{N}}
\newcommand{\Z}{\mathbb{Z}}
\newcommand{\Q}{\mathbb{Q}}
\newcommand{\R}{\mathbb{R}}
\newcommand{\C}{\mathbb{C}}
\newcommand{\F}{\mathbb{F}}
\newcommand{\A}{\mathbb{A}}
\newcommand{\PP}{\mathbb{P}}
\newcommand{\E}{\mathbb{E}}
\newcommand{\eps}{\varepsilon}
\newcommand{\dd}{\,\mathrm{d}}
\begin{document}
\section*{UnsolvedMath ID: 30003098}
\textbf{Source problem code:} catalogue entry
\subsection{Definitions and mathematical statement}
The supplied record combines three questions. Its missing hypotheses and the reported truncation prevent a single unqualified mathematical formulation. The following conditional formulations make standard interpretations explicit; they do not resolve the missing source information.

\begin{enumerate}
\item \textbf{Gorenstein secant cone.} In a candidate formulation over an algebraically closed field \(k\), let
\[
X=\operatorname{Seg}\bigl((\PP^1_k)^5\bigr)\subseteq\PP^{31}_k
\]
be the Segre embedding, and interpret the second secant variety as
\[
\sigma_2(X)=\overline{\bigcup_{x,y\in X,\ x\ne y}\langle x,y\rangle},
\]
where the closure is Zariski closure. If \(I(\sigma_2(X))\) is its homogeneous vanishing ideal, its affine cone has coordinate ring
\[
A=k[z_0,\ldots,z_{31}]/I(\sigma_2(X)).
\]
A Noetherian ring is Gorenstein if every localization at a prime ideal is a Gorenstein local ring; a Noetherian local ring is Gorenstein if it has finite injective dimension as a module over itself. The question, under these conventions, is: \emph{is \(A\) Gorenstein?} The source does not specify the field, its characteristic, the embedding, or the indexing convention for secant varieties. The Segre embedding and the two-point convention above are explicit candidate interpretations.

\item \textbf{Normal polytopes from balls.} Let \(\Lambda\) be a lattice of rank \(d\) and let \(V=\Lambda\otimes_{\Z}\R\). For a bounded ball \(B\subset V\) under a chosen metric and boundary convention, put
\[
P_B=\operatorname{conv}(B\cap\Lambda).
\]
When \(B\cap\Lambda\ne\varnothing\), this is a lattice polytope. Under the ambient-lattice convention, a nonempty lattice polytope \(P\subset V\) is normal if
\[
\forall m\in\Z_{\ge1}\quad
mP\cap\Lambda
=
\{u_1+\cdots+u_m:u_i\in P\cap\Lambda\text{ for }1\le i\le m\},
\]
where \(mP=\{mx:x\in P\}\). For fixed choices of ambient lattice, dimension, and ball convention, the supplied question asks whether \(P_B\) is normal for every ball \(B\). The source does not specify those choices, whether normality uses the ambient lattice or the lattice generated by the polytope, or how balls containing no lattice points are treated. Thus restricting the question to nonempty \(P_B\), or to a particular kind of ball, would require an additional qualification.

\item \textbf{Gradient images and orthogonal subspaces.} A possible algebraic interpretation uses an algebraically closed field \(k\), \(V=k^{n+1}\), and a homogeneous polynomial \(f\in k[x_0,\ldots,x_n]\) of positive degree. Its projective gradient map is the rational map
\[
\gamma_f:\PP(V)\dashrightarrow\PP(V^*),\qquad
[x]\longmapsto[\dd f_x],
\]
with domain
\[
D_f=\{[x]\in\PP(V):\dd f_x\ne0\}.
\]
Interpret a projective linear subspace as \(L=\PP(W)\), where \(W\subseteq V\), and interpret its perpendicular as the projective annihilator
\[
L^\perp=\PP(W^\circ),\qquad
W^\circ=\{\ell\in V^*: \ell(w)=0\text{ for every }w\in W\}.
\]
Under these interpretations, the displayed target is
\[
\PP(W^\circ)\cap\gamma_f\bigl(\PP(W)\cap D_f\bigr)=\varnothing.
\]
For fixed \(f\) and fixed \(r=\dim W\), one possible precise meaning of ``generic'' is that there exists a nonempty Zariski open subset \(U\subseteq\operatorname{Gr}(r,V)\) such that this equality holds for every \(W\in U\). The source leaves unspecified the field and characteristic, the allowed degrees and hypotheses on \(f\), the dimension of \(L\), and the meaning of \(L^\perp\). It also does not say whether ``gradient image'' means this actual rational-map image or its Zariski closure. The positive-degree assumption and all interpretations in this paragraph are conditional, not recovered source hypotheses.
\end{enumerate}

The supplied registry labels the combined record open and its statement \texttt{corrected\_verified}. Its supplied dated assessment nevertheless reports that the identifiable secant-cone question has an affirmative answer for \((\PP^1)^5\), while a truncated second conjecture makes the full extracted claim unclear. These are supplied curation claims, not independently verified mathematical conclusions; the truncated target cannot be reconstructed from this record.
\subsection{Short English statement}
The record asks whether a secant cone is Gorenstein, whether lattice points in every ball have a normal convex hull, and whether a generic projective subspace has gradient image disjoint from its perpendicular. Precise versions depend on missing hypotheses and conventions. Supplied dated curation reports an affirmative answer to the identifiable secant question but retains an open label for the unclear combined record.
\subsection{Sources}
\begin{itemize}
\item Supplied local problem registry: permanent UnsolvedMath ID \(30003098\), source problem code \texttt{catalogue entry}, title ``Gorenstein Secants, Normal Ball Polytopes, and Gradient Images,'' in the Oberwolfach Reports Open Problems collection. The proposed year is recorded as 2016; no individual proposer is supplied.
\item Supplied source citation: \emph{Toric Geometry (2017)}, Oberwolfach Reports, pp.\ 963--1026, DOI \texttt{10.4171/owr/2016/19}. \url{https://doi.org/10.4171/owr/2016/19}. The background describes Oberwolfach Reports as documenting mathematical workshops and their open-problem discussions. The supplied literature-source label for this link is ``OWR 2016 report.''
\item Supplied literature source, labeled ``Michalek--Oeding--Zwiernik'': \url{https://doi.org/10.1093/imrn/rnu056}. No article title, theorem number, or passage identifying its precise application is supplied.
\item Literature assessment checked at \texttt{2026-08-21}: this is supplied dated curation, not independently verified. It reports an affirmative answer to the identifiable secant-cone question for \((\PP^1)^5\), but says that a truncated second conjecture leaves the full extracted claim unclear and therefore normalizes the record to open. The background calls this dated literature triage derived from supplied OWR curation and requests independent verification of time-sensitive status and cited claims. No source was visited for this draft.
\end{itemize}
\end{document}
